Deixe \(\mathbb{M}\) ser o espaço vetorial de matrizes reais \(m \times p\). Para um subespaço vetorial \(S\subseteq \mathbb{M}\), denote por \(\delta(S)\) a dimensão do espaço vetorial gerado por todas as colunas de todas as matrizes em \(S\).

Dizemos que um subespaço vetorial \(T\subseteq \mathbb{M}\) é um \emph{espaço matricial de cobertura} se
\[ \bigcup_{A\in T, A\ne \mathbf{0}} \ker A =\mathbb{R}^p. \]
Tal \(T\) é mínimo se não contém um subespaço vetorial próprio \(S\subset T\) tal que \(S\) também é um espaço matricial de cobertura.
\begin{itemize}
  \item[(a)] Seja \(T\) um espaço matricial de cobertura mínimo e seja \(n=\dim (T)\). Prove que \[ \delta(T)\le \binom{n}{2} \]
  \item[(b)] Prove que para todo inteiro \(n\) podemos encontrar \(m\) e \(p\), e um espaço matricial de cobertura mínimo \(T\) como acima tal que \(\dim T=n\) e \(\delta(T)=\binom{n}{2}\).
\end{itemize}
